\subsection*{自增编号没有给出提交顺序}

现在加入第二项真实业务事实：O1 付款成功后，后来又发生一项经授权的变更。若下游要按业务先后解释它们，除了识别副本属于哪项事件，还要取得事件之间的顺序。

先看自增编号。事务 U1 取得编号 100，暂时没有提交。U2 取得 101 并先提交。扫描器读到 101，把永久游标推进到 101。随后 U1 提交 100。此后只查 \code{id>101}，100 就会永久遗漏。序列取值发生在分配时，事务提交另有次序，回滚也可以留下编号空洞。

因此本章的轮询按待处理状态与到期条件反复扫描，完成的事件才离开该范围。分页或批次可以使用临时游标，一轮结束以后回到完整待办集合。自增编号可以方便地标识记录，事件记录时刻可以解释建行时间，进展覆盖依据则来自尚待处理的状态。

\CNineFigure{ch09-original06}{较早分配的编号可能较晚提交，永久编号游标会越过待办}{fig:ch09-sequence}

即使两项事件已经按业务次序提交，发布也可能改变到达顺序。W1 领取 EV1 后网络阻塞，W2 领取较晚 EV2 并先取得确认，W1 或接管者再发布 EV1。代理的可见次序于是为 EV2、EV1。接管前已经在途的旧副本，还可能在重试以后到达。

RabbitMQ 单个 channel 的发布入队次序有明确作用域。多 channel 的发布序列可以相互交错，重新投递、优先级及多个消费者还会影响实际交付与处理顺序~\cite{rabbitmq-queues}。即使暂时把发送排到一个 channel，数据库中的分配、重新领取与重发也要分别处理。

本章主方案使每项待发事件持续有恢复位置，并允许重复与乱序。要求某订单严格按序可见时，设计者可以在该订单的业务事务内分配序号，限制同一对象的并发发布资格，保持稳定路由，并处理早先未知尝试、重试副本、缺口和消费推进。消费者也按相同序列规则决定何时提交效果。第十章将用具体乱序历史解释消费侧要保存的信息。

\subsection*{重试要保留哪些内容}

EV1 暂时发布失败，当前有资格的发布器把下一次运行时间写入 \code{next\_run\_at}，清除持有者，返回 \code{ready}。等待间隔可以随连续失败增加，并加入随机分散，避免大量工作者同时重连。每次选择有限且可达到的时刻，持续扫描便会再次把 EV1 纳入候选。

PA1、PA2 分别保存尝试身份和通信证据，\code{event\_id}、版本、原快照、目标及其受控语义保持原值。发布器升级、换工作者、重建 channel，都按保存版本读取原事件。目标配置已经不兼容时，事件进入核查，等待授权迁移。核查决定记录原目标怎样映射到同一业务目的，以及哪些旧尝试结果未明，使后续受理证据有确定的解释范围。

\subsection*{已确认以后消息还要留在哪里}

EV1 已 \code{accepted}，本地停止重发。若代理接着因 TTL 删除副本，或队列满时丢掉队首，履约系统可能从未处理过它。这里的 TTL 是消息或队列的存活时长，前章锁键 TTL 则用于结束占用，两者控制的对象与后果不同。

RabbitMQ 的 durable 队列会持久保存队列定义，消息跨重启的保留还依赖消息的持久属性~\cite{rabbitmq-queues}。自动删除、队列 TTL 和管理删除可以结束队列的存续。quorum 队列也支持消息 TTL 和容量限制。容量溢出采用 \code{drop-head} 时可以淘汰旧副本，采用 \code{reject-publish} 则拒绝新发布，发布者继续保存待办。有在途消息时，容量限制允许少量超出，因此容量规划要留出这部分空间~\cite{rabbitmq-quorum-queues}。

消费失败的次数同样影响存续。quorum 队列从 RabbitMQ 4.0 起默认交付上限为 20，超过限制会删除消息或转入死信。4.3 又区分实际失败交付计数与取得计数，部分明确退回不计入交付限制~\cite{rabbitmq-quorum-queues}。要在持续重投期间保留消息，就由部署配置明确它的交付限制与后续去向。

死信表示消息离开原队列，转给另一路由继续处理。quorum 默认死信策略为 \code{at-most-once}，转交过程中可能丢失。较强的转交策略要指定 \code{at-least-once}、相应溢出方式及目标可接受等条件，失败恢复时可能多次转交同一消息，目标据稳定身份处理这些副本~\cite{rabbitmq-quorum-queues}。

本章采用明确的存续前提：目标为持续存在的持久 quorum 队列，发布使用持久消息。关闭消息与队列 TTL、自动删除及自动交付次数丢弃，容量压力选择拒绝新发布，管理流程保持队列和已确认历史。消费者使用手工确认，并在自己的处理依据已经持久保存以后，才允许代理删除消息。交付限制设为无限是这里用于隔离问题的选择，实际系统也要安排持续失败造成的积压与资源占用。

这样，新发布被拒收时，EV1 保留在本地待办中。已经确认的副本按上述保存条件留在代理，直到消费者完成所需的持久转交。业务若采用有限保留或有损淘汰，就增加独立的消费完成证据，以及按原事实核对和补发的规则。上游据这些证据决定何时结束恢复，下游也有据可查地继续未完成效果。

\CNineFigure{ch09-original07}{本地待办、代理已受理副本与下游持久处理之间的继续路径}{fig:ch09-continuation}

\subsection*{毒事件与积压怎样继续处理}

不可路由可以随着拓扑修复而恢复。另一些事件会持续失败，例如保存的版本暂时无人支持、报文超出目标限制，或字段之间存在矛盾。这类反复失败的输入通常称为毒事件。达到重试次数上限时，发布器要保存后续处理归属，让诊断与恢复继续围绕原事实进行。

本章让当前有资格的发布器把 EV1 转成 \code{check}，同时建立持久核查单。核查单有独立 \code{case\_id}，引用 EV1，保存完整原快照的关联、失败观察、负责角色、开放状态和后续决定。发布器在这两项写入共同提交后停止自动重试。人工任务也有可扫描的持久位置，并由负责角色实际领取。

核查者修复路由或部署兼容读取者以后，保存自己的身份、决定编号和处理依据，再把原 EV1 重新安排为 \code{ready}。事件快照保持原值，下一次发布另建 PA。若原事实或内容确实有误，核查者以授权更正事件表达新事实，并保留它与原事件、原核查单的关系。旧副本继续解释原通知，更正事件明确表达后来作出的修正。

积压增加时，可以扩展不同事件的发布并发，并监测最早待办年龄、失败原因、未确认尝试与开放核查单。把最新事件始终放在旧事件之前，会让旧 EV1 永久得不到机会。持续扫描还要提供公平调度，核查队列也要有实际处理能力。

若未受理事件的长期产生速率高于有效处理速率，待办会持续增长。数据库空间、代理容量、工作者和核查资源都有限，系统要限流、拒绝新的受理或增加处理能力。P1 已经成功以后，它的通知责任继续存在。流量控制在后续新购买取得受理以前决定，原 P1 的发布依据保留到有证据结束恢复。

本章保留已受理事件与尝试证据。需要归档时，先把完整语义、事实身份、受理依据和核查决定原子转移到可恢复位置，并保留身份查找关系，再释放活跃扫描空间。旧 \code{ready} 或 \code{check} 可能对应一直没有处理完成的成功事实，归档后的恢复者也要能找到它们。更长期的证据回收将在第十二章结合迟到与遗忘继续讨论。

\subsection*{横向比较：换用 Kafka 改变了哪些恢复位置}

前面的主方案已经闭环：购买库按状态轮询，RabbitMQ 按既定队列条件保留副本，持久核查接住暂时不能发布的工作。现在保持同一业务事件 EV1，只替换消息系统和读取方式，逐项比较哪些恢复依据随之改变。

若把 EV1 改写入 Kafka，还要判断应用恢复后的重发能否自动合并。这里先改变消息的保存模型：发布者向主题 \code{order-events} 追加记录，主题是一组有序分区的集合，每个分区按自己的日志位置保存记录。发布者可以用对象 O1 作为消息键，在固定分区规则与配置下，把该对象的记录选入同一分区。设 EV1 写入分区 2，记录取得该分区的位置，消息内容继续带本地事件号 EV1。日志位置标识这次写入，EV1 关联其业务事实~\cite{kafka-design}。

Kafka 也有名为幂等生产者的能力。其代理按生产者 ID 与发送序号识别协议重试，相关生产配置还限制确认、重试和在途请求数量~\cite{kafka-producer,kafka-design}。Outbox 进程重启以后重新调用一次 \code{send(EV1)}，会成为应用再次提交的记录。同一 EV1 写在消息内容中，代理就按普通消息内容保存，业务事件归并由相应处理协议完成。

换用 Kafka 时，发布器继续保存事件身份与本地恢复关系。消费方按组保存各分区的读取进度，日志则按主题配置保留，某个组推进位置后，其他组还可以读取同一条尚在保留中的记录。第十章从这份组进度推导故障后的重读边界。Kafka 事务可以覆盖它所管理的主题写入和相应消费进度，购买库 R 是另一提交域~\cite{kafka-design}。生产者协议去重减少协议重试造成的副本，应用在恢复时再次提交同一事件，则由稳定 EV1 关联这些记录。

\subsection*{日志订阅改变了哪一段工作}

轮询直接读取待发状态，另一种方式从数据库事务日志读取已提交事件，再转发给代理。事务日志保存数据库变更的提交依据，连接器持续读取并保存自己的进度。这减少了周期扫描的需要，同时把日志保留、连接器进度与重启恢复带入发布链。

Debezium Outbox Event Router 是把这类变更转换成事件消息的组件。其 \code{id} 用于事件身份，\code{aggregateid} 用作消息键，将同一业务对象路由到相应 Kafka 分区~\cite{debezium-outbox}。同一对象使用同一键，可以把它的消息放入同一分区。业务事务先形成不可变事件行，连接器按日志位置恢复读取，消费者再依据自己的处理进度提交效果。各段分别保存恢复位置，才把同一分区的顺序用于实际业务。

该组件对事件表有具体输入规则：预期业务事件以 \code{INSERT} 写入，\code{UPDATE} 按配置警告、报错或停止处理，\code{DELETE} 会被过滤~\cite{debezium-outbox}。本章轮询表会反复更新 \code{phase}、领取和受理状态。接入这个组件时，应把不可变事件行与轮询调度状态分开，让连接器读取前者的新增事件，并为连接器进度和日志保留建立恢复依据。

这两种变体分别改变代理的保存模型与发布器的事件读取路径。本章完整历史继续采用已经推导的轮询发布与 RabbitMQ：数据库保存待发状态，发布器领取事件，代理确认经过 PA 合流到 EV。


\clearpage
\subsection*{从支付成功接到下游业务提交}

把各段放回同一购买：S 绑定 P1/C1，A 共同保存 T1。平台 E 建立 X1，使 10000 分变为 7000 分。R 首次建立 O1、原结果与 EV1，并结束支付任务。PA1 或 PA2 将 EV1 交给指定代理目标，本地保存有效受理证据以后停止继续发布。

安全关系分别落在这些提交点上。R 回滚时，O1、结果和事件共同撤销。R 成功时，事实唯一约束与 C 锁保护同一事件。发布重试保留 EV1 及快照，确认合流只承认有正确归属、路由和持久条件的证据。过期工作者的失败观察保存到自己的 PA，已完成的 EV 保持完成。重复副本可以存在，新的购买则取得自己的 P/C/O 与事件。

持续进展依赖每项待办可发现并公平获调度，数据库恢复可用、持锁事务结束、数据库时钟推进，工作者取得足够连续运行时间。代理与网络最终提供受理机会，原目标拓扑正确且持续可用，至少一次完整确认到达并能持久保存。进入核查的事件由负责者实际作出有依据的决定。已确认副本继续到达消费者，则依前述保存、可用性和消费恢复条件。

即使这些发布条件满足，消费者也可能在业务修改已经提交以后、确认消息以前退出。代理随后再次交付同一 EV1，PA2 形成的另一副本也可能第一次到达该消费者。履约处理现在要根据一份持久记录，判断 O1 的这项效果已经完成，并在退出以后恢复原处理结果。
